\section{Example: Global Magnetic Field Options} \label{ex7}

\W~offers in the namelist \scap{\$B0SCBLOBN} some options,
e.g. shifting, scaling, and symmetry operations of the magnetic field.
It also offers to repeat the set-up periodically. In this example, we make
use of these options in order to
calculate the spontanous synchrotron radiation of a very long
undulator section including phase shifters. But first, we consider a
single cell of this section (Fig. \ref{ex7fig1}). We enter:

\cboxs{\senv \grau \\
python3 ../python/wave\_example\_7.py\\
cp \$WAVE/input/wave\_example\_7.in wave.in\\
wave\\
wplot \\
}

\begin{figure}[htb]
    \centering
    %\framebox(12,8){
    \includegraphics[width=0.8\textwidth]
    {/home/scheer/wav/doc/pp/wave_example_7_single_id.pdf}
    %}
    \caption{Field and 9th harmonic of a single undulator section}
    \label{ex7fig1}
\end{figure}

As the overview shows,
the field is a superposition a planar undulator (namelist \scap{\$ELLIPN}) and
a phase shifter (namelist \scap{\$HALBASY}). If \scap{IBSUPER=1}
in \wie, \W~accepts several field selections . The undulator parameters are
set such, that the 9th harmonic is at \gtis{150\,keV} for a beam energy of
\gtis{17\,GeV}. We now make use of a the periodic field option. Have a
look on the corresponding variables in the namelist
\scap{\$B0SCGLOBN}:\\

\cboxss{IPERIODG=1\\
PERIODG=3.\\
PEROFFG=-1.5\\
}

This means, that the magnetic field is periodically repeated starting from
\scap{X=-1.5} with a period-length of 3 meters.

We set in \wie:

\cboxss{XSTOP=88.5\\
IHTRSMP=10}

and run \W and get the results shown in Fig. \ref{ex7fig3}

\begin{figure}[htb]
    \centering
    %\framebox(12,8){
    \includegraphics[width=0.8\textwidth]
    {/home/scheer/wav/doc/pp/wave_example_7_long.pdf}
    %}
    \caption{Field and 9th harmonic of a long undulator section}
    \label{ex7fig3}
\end{figure}

Passing a long undulator section an electron beam loses a significant
amount of energy. Hence the K-parameter of the down-stream undulators
does not match the beam energy anymore. This results in a broadening
of the spectrum. This can be simulated by setting

\cboxss{IENELOSS=1}

in \wie. We run \W and get the results in Fig. \ref{ex7fig4}.

\begin{figure}[htb]
    \centering
    %\framebox(12,8){
    \includegraphics[width=0.8\textwidth]
    {/home/scheer/wav/doc/pp/wave_example_7_ieneloss.pdf}
    %}
    \caption{Field and 9th harmonic of a long undulator section
    with continous energy loss of the beam}
    \label{ex7fig4}
\end{figure}

To compensate it, we must adjust the field of each undulator
accordingly. This can be achieved with the field option

\cboxss{JBMASK=100}

and by providing a file \gtis{wave.bmask}, which contains a number of
$X$-Intervalls and a scaling factor for the magnetic field for
each intervall. The last \wo gives the total and relative energy loss of the beam due to
the 30 undulators. From this number we calculate the energy loss of a
single undulator cell, and use the relation

\begin{eqnarray*}
   \frac{E_{beam}^2}{1+K^2/2} &=& const
\end{eqnarray*}

to calculate the correction factor for the undulator field, as it is
done in the function \gtis{taper()}, which also writes the file \scap{wave.bmask}

We use this function:

\cboxss{wplot\\
taper()\\
Quit()
}

The default values of this function correspond to this example. Beside the
relative energy loss the starting $K$-value is taken from \wo of the
previous run as well (Look for ''deflection parameter'').

In \wi we set

\cboxss{JBMASK=100}

Please, also read the comments of the corresponding namelist
\scap{B0SCGLOBN} carefully and notice the difference between
\scap{JBMASK} and \scap{IBMASK}.

We run \W~again and see in Fig. \ref{ex7fig5} that the harmonic of the
undulator recovers.

\begin{figure}[htb]
    \centering
    %\framebox(12,8){
    \includegraphics[width=0.8\textwidth]
    {/home/scheer/wav/doc/pp/wave_example_7_taper.pdf}
    %}
    \caption{Field and 9th harmonic of a long undulator section
    with continous energy loss of the beam and corrected undulators}
    \label{ex7fig5}
\end{figure}
